\section{Pseudoindex four}\label{sec:pseudoindex-four}

Watanabe treats this case in \cite[Theorem~2.4]{Watanabe2021}.
We give an intersection-theoretic argument using the factorization of the
universal family established by Kanemitsu. His theorem gives either
rational homogeneity or a factorization of the evaluation map; we exclude
the latter possibility. The relative-degree calculation below reproves
Kanemitsu's Lemma~4.3, and the exclusion of projective-plane fibres
adapts the Betti-number argument in his proof of Theorem~0.2
\cite{Kanemitsu2017}. In dimension six, the remaining quadric-bundle
case is excluded by the orthogonal-bundle calculation given here.

Let $X$ be a smooth connected complex projective Fano sixfold with nef
tangent bundle, Picard number one and pseudoindex four. Let
$\pi:U\to V$ and $e:U\to X$ be the universal family and evaluation map
of a minimal rational component. By Proposition~\ref{prop:minimal-family}, we have
\[
 T_e|_C\simeq\mathcal O_{\mathbb P^1}(-1)^{\oplus2}
 \qquad(C\text{ a }\pi\text{-fibre}).
\]
\begin{proposition}\label{prop:pseudoindex-four}
A smooth complex Fano sixfold with nef tangent bundle, Picard number one
and pseudoindex four is rational homogeneous.
\end{proposition}

\begin{proof}
Kanemitsu's Theorem~3.3 and Proposition~4.1
\cite{Kanemitsu2017} apply in arbitrary dimension; the latter point is
explicit in his Remark~4.2.
They show that either $X$ is rational homogeneous, or there is a diagram
\[
\begin{tikzcd}[column sep=large,row sep=large]
 U \arrow[r,"f"] \arrow[d,"\pi"'] &
 W \arrow[r,"g"] \arrow[d,"q"] & X\\
 V \arrow[r] & Y &
\end{tikzcd}
\]
with $e=g\circ f$, with $f,g,\pi$ smooth $\mathbb P^1$-fibrations,
and with $q$ a smooth elementary Mori contraction, that is, a Mori
contraction of relative Picard number one. We exclude the
second possibility. Here $\dim W=7$. Applied to $f$ and $\pi$,
\cite[Remark~2.3]{Kanemitsu2017} says that all fibres of $q$ are isomorphic to
one member of the list
\[
 \mathbb P^1,\quad\mathbb P^2,\quad\mathbb P^3,\quad
 Q^3,\quad Q^5,\quad K(G_2).
\]
Here $Q^r$ is a smooth quadric of dimension $r$, and $K(G_2)$ is the
five-dimensional rational homogeneous contact manifold of type $G_2$.
In particular, $\dim Y\ge2$.

\smallskip
\noindent\emph{The two relative degrees.}
Write
\[
 H=\frac{c_1(T_X)}4,\qquad h=g^*H,\qquad c=c_1(T_g).
\]
The classes $h,c$ form a basis of $H^2(W,\mathbb Q)$: a smooth
$\mathbb P^1$-fibration has the rational Leray--Hirsch decomposition,
and $H^2(X,\mathbb Q)=\mathbb QH$. These classes are algebraic, so
the same basis describes $N^1(W)_{\mathbb Q}$, the rational divisor
classes modulo numerical equivalence.
Fix a $\pi$-fibre $C$. The map $e|_C$ is birational to its image,
so $f|_C$ is also birational to its image and $h\cdot f(C)=1$.
Set
\[
 a=c_1(T_f)\cdot C,\qquad b=c\cdot f(C).
\]
The relative tangent sequence gives $a+b=-2$. Moreover,
$T_f|_C$ is a line subbundle of $\mathcal O(-1)^{\oplus2}$, whence
$a\le-1$. Both $a$ and $b$ are integers.

We show that $b<0$. The ray contracted by $q$ is generated by $f(C)$.
Since $\rho(W)=2$ and $q$ is elementary, $\rho(Y)=1$. Thus a
positive multiple of $c-bh$ is the pullback
of an ample rational class on $Y$. The sign is positive because
$c-bh$ has degree two on a $g$-fibre, whose numerical class is not
in the contracted ray. By Lemma~\ref{lem:p1-involution},
\[
 g_*c=2,\qquad g_*c^2=0,
\]
and consequently
\[
 g_*(c-bh)^2=-4bH.
\]
Intersecting with any complete-intersection curve in $X$ proves
$b\le0$.

If $b=0$, the square of $c$ has zero integral over $g^{-1}(D)$
for every smooth curve $D\subset X$. Each $g$-fibre has positive
degree against $q^*A$, for an ample class $A$ on $Y$, and hence maps
onto a curve under $q$. Since $q$ is surjective and $\dim Y\ge2$,
two such image curves are distinct. Choose a smooth irreducible
complete-intersection curve $D\subset X$ through the corresponding
two points of $X$. The inverse image $g^{-1}(D)$ is a smooth
irreducible surface, and its image under $q$ contains both of those
distinct curves. Thus this image has dimension two. It follows that
$(q^*A)^2\cdot[g^{-1}(D)]>0$, a contradiction. We have proved
\[
 a=b=-1.
\]
Since $-K_f\cdot C=a=-1$, Kanemitsu's Remark~2.3 now reduces the
possible fibres of $q$ to
$\mathbb P^2$, $Q^3$, and $K(G_2)$.

\smallskip
\noindent\emph{The projective-plane fibres.}
Suppose that $q$ has fibres $\mathbb P^2$, so that $\dim Y=5$.
Relative ample classes and their powers give Leray--Hirsch bases for
both $g$ and $q$. Writing $b_i(Z)=\dim H^i(Z,\mathbb Q)$, we obtain
$b_2(Y)=1$, and Poincar\'e duality
on $Y$ gives
\[
\begin{split}
 b_4(X)+1&=b_4(W)=b_4(Y)+2,\\
 b_6(X)+b_4(X)&=b_6(W)=2b_4(Y)+1.
\end{split}
\]
Hence $b_6(X)=b_4(X)-1$. This contradicts the injectivity of
cup-product by an ample class from $H^4(X,\mathbb Q)$ to
$H^6(X,\mathbb Q)$, which follows from Hard Lefschetz.

\smallskip
\noindent\emph{The remaining fibres.}
We may therefore assume that $\dim Y\le4$. Put
\[
 z=\frac{h+c}{2},\qquad m_i=\int_W h^{7-i}z^i\quad(0\le i\le7).
\]
Since $z\cdot f(C)=0$, the class $z$ is pulled back from
$H^2(Y,\mathbb Q)$. Thus $z^5=0$, and
\[
 m_0=m_5=m_6=m_7=0,\qquad m_1=\int_X H^6>0.
\]
The rational-cohomology involution of $g$ fixes classes pulled back
from $X$ and sends $c$ to $-c$. It sends $z$ to $h-z$ and acts by
$-1$ on top cohomology. For $i=2,4,6$, therefore,
\[
 \sum_{j=0}^i(-1)^j\binom ij m_j=-m_i.
\]
In order, these identities give
\[
 m_2=m_1,\qquad m_1-2m_3+m_4=0,
 \qquad9m_1-20m_3+15m_4=0.
\]
Consequently
\[
 m_1=m_2=5m_4,\qquad m_3=3m_4.
\]
In particular $m_4>0$, so $\dim Y\ge4$. Thus $\dim Y=4$ and
every fibre of $q$ is $Q^3$.

It remains to show that this quadric fibration is incompatible with
$g$. Since $-K_W=4h+c$ has degree three on $f(C)$, write
\[
 -K_q=3h-\lambda z.
\]
The coefficient $\lambda$ is integral: a $g$-fibre has respective
degrees zero and one against $h$ and $z$, and $-K_q$ is an integral
divisor class. We will determine $\lambda$ from the moments above
using the quadric-bundle identity
\[
 q_*(-K_q)^4=0.
\]
To prove this identity, first apply Fischer--Grauert
\cite{FischerGrauert1965}: the smooth projective morphism $q$, all of
whose fibres are $Q^3$, is analytically locally trivial. Its structure
group is $\operatorname{Aut}(Q^3)=\operatorname{PGO}_5(\mathbb C)$.
In odd rank this group is isomorphic to $\operatorname{SO}_5$.
Explicitly, fix a nondegenerate symmetric $5\times5$ matrix $J$ defining
the quadric. If $A$ represents a projective automorphism and
$A^{\mathsf T}JA=\mu J$, then
\[
 \widetilde A=\frac{\mu^2}{\det A}A
\]
has determinant one and preserves $J$. This expression is unchanged
when $A$ is multiplied by a scalar; it is the unique special
orthogonal lift and respects products. The transition functions thus
define a rank-five holomorphic vector bundle $E$ equipped with a
nondegenerate symmetric form with values in $\mathcal O_Y$ and
trivial determinant. GAGA algebraizes this bundle and its form
\cite[Theorems~2--3 and Propositions~15, 18]{Serre1956}.
The resulting analytic identification of $W$ with the relative quadric
is algebraic as well.
Accordingly, in the quotient convention for projectivization,
$W\subset\mathbb P_Y(E)$ is the quadric of class $2\zeta$, where
$\zeta=c_1(\mathcal O_{\mathbb P(E)}(1))$. Self-duality identifies
the line and quotient descriptions of this quadric. If
$p:\mathbb P_Y(E)\to Y$ is the projection, adjunction gives
\[
 -K_q=3\zeta|_W,
 \qquad q_*(-K_q)^4=2\cdot3^4p_*\zeta^5
 =2\cdot3^4c_1(E)=0.
\]

Now the projection formula and $z^5=0$ give
\[
 0=\int_W(-K_q)^4z^3=3^4m_3-4\cdot3^3\lambda m_4.
\]
Since $m_3=3m_4$ and $m_4>0$, this forces $\lambda=9/4$,
contrary to integrality. This excludes the factorization case and
proves the proposition.
\end{proof}
